% =====================================================================
%  Lezione 8 — Teoria del matching e applicazioni
%  Materiale: Lesson_8_Matching_Theory.pdf (nessun notebook di laboratorio:
%             la parte LAB implementa e verifica gli algoritmi delle slide)
% =====================================================================
\chapter[Teoria del matching]{Teoria del matching\\e applicazioni}

\parteteoria

Un problema di \kw{matching} (abbinamento) nasce quando degli agenti hanno preferenze su
altri agenti o su oggetti, e bisogna decidere chi va con chi. Gli esempi delle slide sono
molto concreti: scambiare case o reni, distribuire $n$ oggetti fra $n$ persone, assegnare
specializzandi agli ospedali e studenti alle scuole, formare coppie, assegnare i revisori
agli articoli o le persone ai compiti. Il percorso:
\begin{enumerate}[label=\arabic*.]
  \item preferenze da \kw{un lato solo}: la dittatura seriale e lo scambio con il
        \kw{Top Trading Cycle};
  \item preferenze da \kw{entrambi i lati}: il matrimonio stabile e l'algoritmo di
        \kw{Gale--Shapley}, con le sue proprietà e i suoi limiti;
  \item il matching come \kw{ottimizzazione}: flusso massimo e flusso a costo minimo.
\end{enumerate}
Torna un tema della lezione 7: una regola di decisione collettiva va giudicata per
efficienza, equità e \en{strategyproofness}, e di solito non si possono avere tutte.

% ---------------------------------------------------------------------
\section{Il panorama dei problemi di matching}

La lezione tratta il \kw{matching bipartito}: due insiemi disgiunti (persone e oggetti,
studenti e scuole, uomini e donne) e ogni abbinamento collega un elemento del primo con uno
del secondo.

\begin{definizione}[Matching]
Dati due insiemi disgiunti $U$ e $V$, un \kw{matching} è un insieme di coppie
$(u, v) \in U \times V$ in cui ogni elemento compare al più una volta. È \kw{perfetto} se
ogni elemento di $U$ e di $V$ compare esattamente una volta (serve $|U| = |V|$).
\end{definizione}

Tre famiglie, che danno un significato diverso a ``buon matching'':
\begin{itemize}
  \item \kw{preferenze da un lato}: solo un lato ha preferenze (le persone sugli oggetti).
        Si cercano efficienza, equità e incentivi a dire il vero;
  \item \kw{preferenze da due lati}: entrambi i lati hanno preferenze sull'altro. La
        proprietà centrale diventa la \kw{stabilità}: nessuna coppia vuole lasciare
        l'abbinamento assegnato;
  \item \kw{ottimizzazione}: si vuole massimizzare o minimizzare una funzione globale
        (numero di coppie, costo totale), anche dove la stabilità non si può ottenere.
\end{itemize}

% ---------------------------------------------------------------------
\section{Preferenze da un lato: la dittatura seriale}

Ci sono $n$ persone e $n$ oggetti; ogni persona ha un ordinamento stretto sugli oggetti e
ne deve ricevere esattamente uno. L'esempio delle slide, con le liste dalla scelta migliore
alla peggiore:
\[
  A = [X, Y, Z], \qquad B = [Z, X, Y], \qquad C = [Z, Y, X].
\]

\begin{definizione}[Dittatura seriale]
Si fissa un ordine degli agenti. Seguendo l'ordine, ogni agente prende l'oggetto che
preferisce fra quelli ancora liberi. Se l'ordine è estratto a caso si parla di
\kw{dittatura seriale casuale} (\en{random serial dictatorship}).
\end{definizione}

\begin{esempio}[l'ordine conta]
Con l'ordine $A, B, C$: $A$ prende $X$, $B$ prende $Z$, a $C$ resta $Y$. È il matching delle
slide, $\{(A,X), (B,Z), (C,Y)\}$: $A$ e $B$ hanno la prima scelta, $C$ la seconda. Con
l'ordine $C, B, A$: $C$ prende $Z$, $B$ prende $X$, ad $A$ resta $Y$. Cambiando l'ordine
cambia chi resta deluso.
\end{esempio}

Le proprietà elencate nelle slide:
\begin{itemize}
  \item \kw{Pareto efficiente}: non si può migliorare qualcuno senza peggiorare qualcun
        altro. Il primo agente ha il suo oggetto preferito; per dare a un agente qualcosa di
        meglio bisognerebbe toglierlo a uno che l'ha scelto prima di lui, e che quindi
        ci perde;
  \item \kw{strategyproof}: mentire non conviene. Quando arriva il proprio turno si prende
        il migliore fra quelli rimasti, e gli oggetti rimasti non dipendono da ciò che
        si dichiara;
  \item \kw{non massimizza il benessere sociale}: con utilità cardinali può esistere
        un'allocazione con somma più alta;
  \item \kw{equa?} Con l'ordine casuale tutti gli agenti sono trattati allo stesso modo
        \emph{prima} dell'estrazione: è equa in valore atteso, non a posteriori.
\end{itemize}

\begin{esempio}[non massimizza il benessere]
Due agenti e due oggetti. $A$ dà valore 10 a $X$ e 9 a $Y$; $B$ dà valore 10 a $X$ e 0 a
$Y$. Con l'ordine $A, B$: $A$ prende $X$, $B$ resta con $Y$, benessere $10 + 0 = 10$.
L'allocazione $(A, Y), (B, X)$ dà $9 + 10 = 19$. La dittatura seriale usa solo gli
ordinamenti: non vede che per $A$ la differenza fra $X$ e $Y$ è piccola e per $B$ enorme.
\end{esempio}

\subsection{Estensioni}

\begin{itemize}
  \item \kw{Numero diverso di agenti e oggetti} ($n \ne m$): funziona senza modifiche. Con
        più oggetti che agenti alcuni oggetti restano liberi; con più agenti gli ultimi
        restano senza.
  \item \kw{Capacità}: ogni agente può ricevere più oggetti (e ogni oggetto può avere più
        copie). Due varianti:
        \begin{itemize}
          \item \kw{dittatura seriale fino alla capacità}: nell'ordine, ogni agente prende
                in un colpo solo tanti oggetti quanta è la sua capacità. Strategyproof ma
                iniqua: il primo prende tutto il meglio;
          \item \kw{round robin}: a ogni turno ogni agente prende \emph{un} oggetto, e si
                ripete finché le capacità sono piene. Più equa, ma non strategyproof.
        \end{itemize}
\end{itemize}

\begin{esempio}[manipolare il round robin]
Due agenti, quattro oggetti, due a testa; l'ordine è $1, 2, 1, 2$. Le preferenze vere sono
$1: a \succ b \succ c \succ d$ e $2: b \succ c \succ a \succ d$, con utilità additive.
Dichiarando il vero: $1$ prende $a$, $2$ prende $b$, $1$ prende $c$, $2$ prende $d$;
l'agente 1 ha $\{a, c\}$. Se invece l'agente 1 dichiara $b \succ a \succ c \succ d$: prende
$b$, l'agente 2 (che voleva $b$) prende $c$, l'agente 1 prende $a$. Ora ha $\{a, b\}$, che
preferisce. Ha senso prendere subito l'oggetto conteso e lasciare per dopo quello che
nessun altro vuole.
\end{esempio}

La distribuzione di oggetti \emph{indivisibili} fra agenti è studiata nell'AGT sotto il
nome di \kw{fair division of indivisible goods}, uno degli argomenti ``in arrivo'' citati
nella lezione 1.

% ---------------------------------------------------------------------
\section{Scambi: l'algoritmo Top Trading Cycle}

Situazione diversa: ci sono ancora $n$ agenti e $n$ oggetti, ma ogni agente \kw{possiede
già} un oggetto e ha preferenze strette su tutti. Si vuole scambiare gli oggetti in modo che
ognuno stia almeno come prima. Scenari tipici: programmi di edilizia sociale, scambio di
stanze, scambio di reni fra coppie donatore--paziente incompatibili. Il modello si chiama
\kw{housing market} (Shapley e Scarf, 1974).

\begin{definizione}[Top Trading Cycle, TTC]
Algoritmo attribuito a David Gale e pubblicato da Shapley e Scarf:
\begin{enumerate}
  \item si costruisce un grafo orientato: ogni agente punta al suo oggetto preferito fra
        quelli rimasti, ogni oggetto punta al suo proprietario;
  \item il grafo ha almeno un ciclo: se ne trova uno e ogni agente del ciclo riceve
        l'oggetto a cui punta;
  \item si tolgono gli agenti e gli oggetti del ciclo e si ripete.
\end{enumerate}
\end{definizione}

\begin{osservazione}[perché un ciclo c'è sempre]
Nel grafo ogni nodo ha esattamente un arco uscente. Partendo da un nodo qualunque e
seguendo gli archi, prima o poi si torna su un nodo già visto, perché i nodi sono finiti:
quello è un ciclo. Il ciclo può anche essere lungo uno, cioè un agente punta al proprio
oggetto: in quel caso se lo tiene.
\end{osservazione}

\begin{esempio}[il TTC delle slide]
Quattro agenti; fra parentesi l'oggetto posseduto.
\[
\begin{array}{ll}
  A\ (i_1):\ [i_3, i_2, i_4, i_1] & \qquad B\ (i_2):\ [i_4, i_1, i_2, i_3] \\
  C\ (i_3):\ [i_1, i_4, i_3, i_2] & \qquad D\ (i_4):\ [i_3, i_2, i_1, i_4]
\end{array}
\]
\textbf{Primo giro.} $A \to i_3 \to C \to i_1 \to A$ è un ciclo: $A$ riceve $i_3$ e $C$
riceve $i_1$. Anche $B$ e $D$ puntano a $i_4$ e a $i_3$, ma non chiudono un ciclo.
\begin{center}
\begin{tikzpicture}[x=1cm, y=1cm, font=\small,
    ag/.style={circle, draw=blu, thick, minimum size=6.5mm, inner sep=0pt},
    og/.style={circle, draw=unical, thick, minimum size=6.5mm, inner sep=0pt},
    fr/.style={-{Stealth[length=5pt]}, thick}]
  \foreach \n/\y in {A/3, B/2, C/1, D/0} \node[ag] (\n) at (0,\y) {\n};
  \foreach \n/\y in {1/3, 2/2, 3/1, 4/0} \node[og] (i\n) at (3.2,\y) {$i_\n$};
  \draw[fr, unical, very thick] (A) -- (i3);
  \draw[fr, unical, very thick] (C) -- (i1);
  \draw[fr] (B) -- (i4);
  \draw[fr] (D) -- (i3);
  \draw[fr, grigio, dashed] (i1) to[bend right=25] (A);
  \draw[fr, grigio, dashed] (i2) to[bend right=25] (B);
  \draw[fr, grigio, dashed] (i3) to[bend right=25] (C);
  \draw[fr, grigio, dashed] (i4) to[bend right=25] (D);
  \node[anchor=west, text width=5.2cm, font=\footnotesize] at (4.3,1.5)
    {agente $\to$ oggetto preferito (pieno);\\ oggetto $\to$ proprietario (tratteggiato).\\
     In rosso il ciclo $A \to i_3 \to C \to i_1 \to A$.};
\end{tikzpicture}
\end{center}
\textbf{Secondo giro.} Restano $B, D$ e $i_2, i_4$. $B$ punta a $i_4$ (di $D$), $D$ non
può più avere $i_3$ e punta a $i_2$ (di $B$): ciclo $B \to i_4 \to D \to i_2 \to B$.
Risultato: $(A, i_3), (C, i_1), (B, i_4), (D, i_2)$. Tre agenti hanno la prima scelta, $D$
la seconda; tutti stanno meglio di prima.
\end{esempio}

Le proprietà del TTC:
\begin{itemize}
  \item \kw{individualmente razionale}: nessuno riceve un oggetto peggiore del proprio,
        perché il proprio oggetto resta disponibile finché il suo proprietario è in gioco;
  \item \kw{stabile}, nel senso del \kw{core}: nessun gruppo di agenti può fare meglio
        scambiandosi fra loro solo i propri oggetti. Con preferenze strette il risultato è
        l'\kw{unica allocazione nel core} (Roth e Postlewaite, 1977);
  \item \kw{strategyproof} (Roth, 1982): dichiarare preferenze false non fa ottenere un
        oggetto migliore;
  \item \kw{Pareto efficiente}: segue dal core, perché la coalizione di tutti gli agenti
        non può migliorare.
\end{itemize}

\begin{attenzione}[cosa garantisce davvero]
La slide di presentazione dice che lo scopo è ``che ogni agente ottenga un oggetto
migliore''. Non è sempre possibile: se tutti preferiscono il proprio oggetto, il TTC li
lascia dove sono (tutti cicli di lunghezza uno). La garanzia è più debole e più precisa:
\emph{nessuno peggiora}, e nessun gruppo può fare meglio da solo. Anche ``nessun agente
vuole deviare'' va letto così: è la stabilità del core, non quella delle coppie che vedremo
nel matrimonio stabile.
\end{attenzione}

\begin{extra}[lo scambio di reni]
Una coppia paziente--donatore incompatibile ``possiede'' un rene che non può usare. Il TTC
troverebbe cicli di scambio, ma nella pratica ogni ciclo richiede operazioni simultanee,
quindi si limitano i cicli a lunghezza 2 o 3. Con questo vincolo il problema di massimizzare
i trapianti diventa NP-difficile e si risolve con la programmazione lineare intera. Roth e
Shapley hanno ricevuto il Nobel per l'economia nel 2012 proprio per la teoria delle
allocazioni stabili e la sua applicazione a mercati come questo.
\end{extra}

% ---------------------------------------------------------------------
\section{Preferenze da due lati: il matrimonio stabile}

Ora entrambi i lati hanno preferenze. La formulazione classica: $n$ uomini e $n$ donne, ogni
agente ha un ordinamento stretto e completo sugli agenti dell'altro lato. Si vuole un
matching perfetto \kw{stabile}.

\begin{definizione}[Coppia bloccante e matching stabile]
Dato un matching $\mu$, una coppia $(m, w)$ che non è in $\mu$ è \kw{bloccante} se $m$
preferisce $w$ alla sua partner $\mu(m)$ \emph{e} $w$ preferisce $m$ al suo partner
$\mu(w)$. Un matching è \kw{stabile} se non ha coppie bloccanti.
\end{definizione}

In parole: nessun uomo e nessuna donna vogliono ``divorziare'' dai partner assegnati per
mettersi insieme. Per bloccare serve che \emph{entrambi} ci guadagnino: se solo uno dei due
vuole cambiare, la coppia non si forma.

\begin{attenzione}[stabilità non è efficienza]
La slide descrive la stabilità anche come ``non si può riassegnare una coppia per
migliorarla senza peggiorare qualcun altro''. Questa è la \emph{Pareto efficienza}, un
concetto diverso. Con preferenze strette ogni matching stabile è anche Pareto efficiente:
se $\mu'$ migliorasse qualcuno, per esempio $m$ con $\mu'(m) = w$, senza peggiorare nessuno,
anche $w$ avrebbe cambiato partner, quindi starebbe strettamente meglio con $m$, e $(m, w)$
bloccherebbe $\mu$. Il contrario è falso: ci sono matching Pareto efficienti instabili. La
definizione da usare è quella con le coppie bloccanti.
\end{attenzione}

\section{L'algoritmo di Gale--Shapley}

\begin{teorema}[Gale e Shapley, 1962]
Per qualunque profilo di preferenze di $n$ uomini e $n$ donne esiste un matching stabile.
\end{teorema}

La dimostrazione è costruttiva: l'algoritmo di \kw{accettazione differita}
(\en{deferred acceptance}) ne trova uno.
\begin{itemize}
  \item A ogni turno, ogni uomo non fidanzato fa la proposta alla donna che preferisce fra
        quelle a cui non l'ha ancora fatta.
  \item A ogni turno, ogni donna sceglie il preferito fra le proposte ricevute e il
        fidanzato attuale; rifiuta gli altri.
\end{itemize}
``Differita'' perché nessuna accettazione è definitiva: una donna tiene il fidanzato solo
finché non arriva qualcuno che preferisce. L'algoritmo finisce quando tutti sono fidanzati.

\begin{lstlisting}[style=pseudo]
tutti liberi
while esiste un uomo libero m con donne ancora da provare:
    w = la prima donna nella lista di m a cui non ha ancora proposto
    if w e' libera: (m, w) si fidanzano
    else if w preferisce m al fidanzato m': (m, w) si fidanzano, m' torna libero
    else: w rifiuta m
return i fidanzamenti
\end{lstlisting}

\begin{esempio}[il Gale--Shapley delle slide]
\[
\begin{array}{lll}
  A = [X, Y, Z] & \qquad & X = [A, B, C] \\
  B = [Z, X, Y] & \qquad & Y = [B, C, A] \\
  C = [Z, Y, X] & \qquad & Z = [A, C, B]
\end{array}
\]
Proposte degli uomini, nell'ordine delle slide:
\begin{center}
\renewcommand{\arraystretch}{1.15}
\begin{tabular}{rll}
\toprule
Passo & Proposta & Esito \\
\midrule
1 & $A \to X$ & $X$ è libera: $(A, X)$ \\
2 & $B \to Z$ & $Z$ è libera: $(B, Z)$ \\
3 & $C \to Z$ & $Z$ preferisce $C$ a $B$: $(C, Z)$, $B$ torna libero \\
4 & $B \to X$ & $X$ preferisce $A$ a $B$: rifiutato \\
5 & $B \to Y$ & $Y$ è libera: $(B, Y)$ \\
\bottomrule
\end{tabular}
\end{center}
Risultato $\{(A, X), (B, Y), (C, Z)\}$, con 5 proposte. Controllo della stabilità: $A$ e
$C$ hanno la prima scelta (per $C$, $Z$ è in cima) e non bloccano nulla; $B$ preferirebbe
$Z$ e $X$, ma $Z$ preferisce $C$ e $X$ preferisce $A$. Nessuna coppia bloccante. In questo
profilo il matching stabile è anche l'unico: lo si verifica provando tutte le $3! = 6$
assegnazioni, e lo stesso risultato esce facendo proporre le donne.
\end{esempio}

\begin{teorema}[correttezza dell'accettazione differita]
L'algoritmo termina dopo al più $n^2$ proposte e restituisce un matching perfetto stabile.
\end{teorema}
\paragraph{Terminazione.} Ogni uomo propone a ogni donna al più una volta: al più $n$
proposte a testa, $n^2$ in tutto. Le slide dicono ``$O(n^2)$ passi''.

\paragraph{Matching perfetto.} Una donna che ha ricevuto una proposta resta fidanzata fino
alla fine (cambia solo partner, e sempre in meglio). Se alla fine un uomo $m$ fosse libero,
avrebbe proposto a tutte le $n$ donne, che quindi sarebbero tutte fidanzate: $n$ donne con
al più $n - 1$ uomini, impossibile.

\paragraph{Stabilità.} Supponiamo che $(m, w)$ sia bloccante: $m$ preferisce $w$ alla
partner finale $\mu(m)$, e $w$ preferisce $m$ a $\mu(w)$. Siccome $m$ propone in ordine di
preferenza, ha proposto a $w$ prima che a $\mu(m)$. Quindi $w$ l'ha rifiutato, subito o più
tardi, per qualcuno che preferiva; e da allora i suoi fidanzati sono solo migliorati. Il
partner finale $\mu(w)$ è quindi preferito a $m$: contraddizione. (È l'argomento delle
slide: ``o $m$ non ha proposto a $w$, o $w$ l'ha rifiutato per uno peggiore, ed entrambe le
cose sono impossibili''.)

\begin{osservazione}[l'ordine delle proposte non conta]
Le slide descrivono turni in cui tutti gli uomini liberi propongono insieme; l'esempio fa
una proposta alla volta. Il risultato è lo stesso: qualunque ordine si scelga per le
proposte, l'accettazione differita con gli uomini che propongono restituisce sempre lo
stesso matching, quello ottimo per gli uomini del prossimo paragrafo.
\end{osservazione}

% ---------------------------------------------------------------------
\section{Chi propone vince: ottimalità e strategie}

Un profilo può avere più matching stabili. Ci si chiede quale restituisce l'algoritmo e chi
ne trae vantaggio.

\begin{definizione}[partner stabile, M-ottimale, W-ottimale]
$w$ è un \kw{partner stabile} di $m$ se esiste un matching stabile che li mette insieme. Un
matching stabile è \kw{M-ottimale} se ogni uomo lo preferisce almeno quanto ogni altro
matching stabile, cioè ogni uomo ha la migliore fra le sue partner stabili. È
\kw{W-ottimale} se vale lo stesso per ogni donna. Simmetricamente si definiscono
\kw{M-pessimale} e \kw{W-pessimale}.
\end{definizione}

Non è ovvio che un matching M-ottimale esista: potrebbe darsi che le migliori partner
stabili di due uomini coincidano. Il teorema seguente lo dimostra mostrando che
l'algoritmo lo costruisce.

\begin{teorema}[ottimalità di chi propone]
L'accettazione differita con proposte degli uomini restituisce il matching stabile
M-ottimale, che è anche W-pessimale. Facendo proporre le donne si ottiene il W-ottimale,
che è M-pessimale.
\end{teorema}

\paragraph{M-ottimalità.} Mostriamo che nessun uomo viene mai rifiutato da una sua partner
stabile. Supponiamo il contrario e prendiamo il \emph{primo} rifiuto di questo tipo durante
l'esecuzione: $w$ rifiuta $m$, partner stabile di $m$, perché ha (o riceve) $m'$ che
preferisce. Sia $\mu'$ un matching stabile con $(m, w)$. In $\mu'$, $m'$ è con un'altra
donna $w'$. Quando $m'$ ha proposto a $w$ non era ancora stato rifiutato da nessuna partner
stabile (questo è il primo rifiuto del genere), quindi non era ancora stato rifiutato da
$w'$: siccome propone in ordine, preferisce $w$ a $w'$. Ma allora in $\mu'$ la coppia
$(m', w)$ è bloccante: $m'$ preferisce $w$ a $w'$, e $w$ preferisce $m'$ a $m$. Assurdo.
Quindi ogni uomo finisce con la prima partner stabile della sua lista, la migliore.

\paragraph{W-pessimalità.} Sia $\mu$ il risultato e $(m, w) \in \mu$. Se esistesse un
matching stabile $\mu'$ in cui $w$ sta peggio, con $m'$ al posto di $m$, allora in $\mu'$
l'uomo $m$ sta con $w' \ne w$. Per la M-ottimalità $w$ è la migliore partner stabile di
$m$, quindi $m$ preferisce $w$ a $w'$; e $w$ preferisce $m$ a $m'$. La coppia $(m, w)$
bloccherebbe $\mu'$: contraddizione.

\begin{attenzione}[le dimostrazioni nelle slide]
Le due dimostrazioni delle slide vanno nella direzione giusta ma saltano un passaggio. Per
l'ottimalità dicono che, se $m$ non ottiene la donna preferita $w$, è perché $w$ l'ha
rifiutato per $m'$, e quindi un matching con $(m, w)$ non sarebbe stabile. Il passaggio
mancante è perché $m'$ preferisca $w$ alla partner che avrebbe in quell'altro matching:
serve l'argomento del \emph{primo} rifiuto qui sopra. E l'enunciato giusto confronta $m$
con le sole partner \emph{stabili}: la donna in cima alla sua lista può benissimo non
essere raggiungibile. Allo stesso modo, nella W-pessimalità la slide scrive che ``$m$
preferisce $w$ a ogni altra donna della sua lista''; il vero è ``a ogni altra partner
stabile''.
\end{attenzione}

\subsection{Nessun meccanismo stabile è strategyproof}

Un \kw{meccanismo di matching} riceve le preferenze dichiarate e restituisce un matching. È
\kw{stabile} se il matching è stabile rispetto alle preferenze dichiarate; è
\kw{strategyproof} se a nessun agente conviene dichiarare il falso.

\begin{teorema}[Roth, 1982]
Non esiste un meccanismo di matching che sia sempre stabile e strategyproof.
\end{teorema}
\paragraph{Dimostrazione (dalle slide).} Basta un profilo in cui ogni scelta stabile
lascia a qualcuno un incentivo a mentire:
\[
\begin{array}{lll}
  m_1: w_1, w_2, w_3 & \qquad & w_1: m_2, m_1, m_3 \\
  m_2: w_2, w_1, w_3 & \qquad & w_2: m_1, m_2, m_3 \\
  m_3: w_1, w_2, w_3 & \qquad & w_3: m_1, m_2, m_3
\end{array}
\]
Gli unici matching stabili sono due:
\[
  \mu_1 = \{(m_1, w_1), (m_2, w_2), (m_3, w_3)\}, \qquad
  \mu_2 = \{(m_1, w_2), (m_2, w_1), (m_3, w_3)\}.
\]
$\mu_1$ è M-ottimale (ogni uomo ha la prima scelta possibile), $\mu_2$ è W-ottimale.
\begin{itemize}
  \item Se il meccanismo sceglie $\mu_1$, la donna $w_1$ può dichiarare
        $w_1: m_2, m_3, m_1$. Con questa lista la coppia $(m_3, w_1)$ blocca $\mu_1$, e
        l'unico matching stabile diventa $\mu_2$: $w_1$ ottiene $m_2$, la sua vera prima
        scelta.
  \item Se il meccanismo sceglie $\mu_2$, l'uomo $m_1$ può dichiarare $m_1: w_1, w_3, w_2$.
        Ora $(m_1, w_3)$ blocca $\mu_2$ e l'unico matching stabile è $\mu_1$: $m_1$ ottiene
        $w_1$, la sua vera prima scelta.
\end{itemize}
In entrambi i casi un meccanismo stabile deve restituire il matching favorevole a chi ha
mentito. \hfill$\square$

Tutto verificato enumerando le $3! = 6$ assegnazioni per ciascuno dei tre profili (quello
vero e i due manipolati).

\begin{osservazione}[chi può mentire]
Il teorema non dice che tutti possono manipolare. Un risultato di Dubins e Freedman (1981)
e di Roth (1982) mostra che l'accettazione differita è strategyproof \emph{per chi
propone}: nessun uomo ottiene di più mentendo quando propongono gli uomini. Sono le donne,
il lato che riceve le proposte, a poter manipolare, come $w_1$ nell'esempio. È per questo
che nei sistemi reali di ammissione si fanno proporre gli studenti: sono loro a cui si
chiede di dichiarare il vero.
\end{osservazione}

% ---------------------------------------------------------------------
\section{Estensioni del matrimonio stabile}

\begin{itemize}
  \item \kw{$n$ uomini e $m$ donne con $n \ne m$}: l'accettazione differita funziona
        ancora. O si cambia il criterio di arresto (un uomo rifiutato da tutte resta
        solo), o si aggiungono individui fittizi, messi in fondo a tutte le liste, per avere
        $n = m$.
  \item \kw{Liste incomplete}: un uomo non propone alle donne che considera inaccettabili e
        una donna non accetta uomini inaccettabili. Ora alcuni restano soli, e una coppia
        blocca anche quando entrambi preferiscono stare insieme piuttosto che soli.
  \item \kw{Preferenze con pareggi}: si rompono i pareggi a caso e si esegue l'algoritmo.
\end{itemize}

\begin{extra}[pareggi e liste incomplete insieme]
Rompendo i pareggi a caso si ottiene un matching \emph{debolmente} stabile (nessuna coppia
in cui entrambi stanno strettamente meglio), ma diverse rotture possono dare matching di
dimensioni diverse. Trovare il più grande matching debolmente stabile con pareggi e liste
incomplete è NP-difficile (Manlove et al., 2002). Con liste incomplete e preferenze strette,
invece, tutti i matching stabili lasciano soli gli stessi agenti (\en{teorema degli
ospedali rurali}).
\end{extra}

\subsection{Matching stabile con capacità}

Studenti e scuole: gli studenti hanno preferenze sulle scuole, le scuole sugli studenti, e
la scuola $i$ ha al più $c_i$ posti. Con $n$ studenti e $m$ scuole di capacità
$c_1, \dots, c_m$ si usa l'accettazione differita in cui la scuola $i$ tiene fino a $c_i$
studenti, i migliori secondo le sue preferenze. Una coppia studente--scuola è bloccante se
lo studente preferisce quella scuola alla propria e la scuola ha un posto libero oppure
preferisce lo studente al peggiore fra quelli ammessi. È il modello degli ospedali e degli
specializzandi (\en{hospitals/residents}) e delle ammissioni scolastiche.

\begin{esempio}[le capacità delle slide]
Cinque studenti, tutti con $[X, Y]$. Le scuole: $X = [A, E, B, C, D]$ con $c_X = 2$,
$Y = [E, A, B, C, D]$ con $c_Y = 3$. Proposte degli studenti:
\begin{center}
\renewcommand{\arraystretch}{1.15}
\begin{tabular}{rll}
\toprule
Passo & Proposta & Esito \\
\midrule
1--2 & $A \to X$, $B \to X$ & $X$ tiene $\{A, B\}$: è piena \\
3 & $C \to X$ & rifiutato: $X$ preferisce $A$ e $B$ \\
4 & $D \to X$ & rifiutato: $X$ preferisce $A$ e $B$ \\
5 & $E \to X$ & $X$ preferisce $E$ a $B$: tiene $\{A, E\}$, $B$ è rifiutato \\
6--8 & $B, C, D \to Y$ & $Y$ ha tre posti: tiene $\{B, C, D\}$ \\
\bottomrule
\end{tabular}
\end{center}
Risultato: $X = \{A, E\}$, $Y = \{B, C, D\}$. È stabile: $B$, $C$, $D$ vorrebbero $X$, ma
$X$ è piena con due studenti che preferisce a ciascuno di loro. Nella slide del passo 4 il
commento dice ancora ``$A, B$ preferiti a $C$'': si intende $D$.
\end{esempio}

% ---------------------------------------------------------------------
\section{Oltre la stabilità: matching e flussi}

A volte la stabilità non si può ottenere, oppure non è ciò che interessa: si vuole
massimizzare il numero di coppie o minimizzare un costo totale. Molti di questi problemi si
riducono a un problema di \kw{flusso su rete}, per cui esistono algoritmi polinomiali.

\begin{attenzione}[non tutto si riduce a un flusso]
Alcuni problemi di matching sono NP-difficili, per esempio il \kw{generalized assignment}
(compiti con dimensioni diverse da assegnare ad agenti con capacità, a costo minimo). Gli
algoritmi di flusso non li risolvono esattamente, ma si possono usare per ottenere soluzioni
approssimate, per esempio risolvendo il rilassamento e poi arrotondando.
\end{attenzione}

\subsection{Reti di flusso e flusso massimo}

\begin{definizione}[Rete di flusso]
Una \kw{rete di flusso} è un grafo orientato in cui ogni arco $e$ ha una \kw{capacità}
$c(e) \ge 0$, con una \kw{sorgente} $s$ e un \kw{pozzo} $t$. Un flusso assegna a ogni arco
una quantità $f(e)$ con:
\begin{enumerate}
  \item \kw{vincolo di capacità}: $0 \le f(e) \le c(e)$;
  \item \kw{conservazione}: in ogni nodo diverso da $s$ e $t$ entra quanto esce,
        \[ \sum_{e \text{ entra in } v} f(e) = \sum_{e \text{ esce da } v} f(e). \]
\end{enumerate}
Il \kw{valore} del flusso è quanto esce da $s$ (uguale a quanto entra in $t$).
\end{definizione}

Si pensi ad acqua in una rete di tubi, a dati fra router, a merci lungo una catena di
fornitura.

\begin{center}
\begin{tikzpicture}[x=1.45cm, y=1.15cm, font=\small,
    nd/.style={circle, draw, thick, minimum size=6.5mm, inner sep=0pt},
    st/.style={circle, fill=blu, text=white, minimum size=6.5mm, inner sep=0pt},
    fr/.style={-{Stealth[length=5pt]}, blu!70},
    lb/.style={fill=white, inner sep=1pt, font=\footnotesize}]
  \node[st] (s) at (0,0) {$s$};
  \node[nd] (a) at (1.6,1.3) {$a$};
  \node[nd] (b) at (1.6,0) {$b$};
  \node[nd] (c) at (1.6,-1.3) {$c$};
  \node[nd] (d) at (4,1.3) {$d$};
  \node[nd] (e) at (4,0) {$e$};
  \node[nd] (f) at (4,-1.3) {$f$};
  \node[st] (t) at (5.6,0) {$t$};
  \foreach \u/\v/\c/\p in {s/a/10/0.5, s/b/5/0.5, s/c/15/0.5, a/b/4/0.5, b/c/4/0.5,
      a/d/9/0.5, a/e/15/0.3, b/e/8/0.5, b/f/6/0.3, c/f/13/0.5, d/e/15/0.5, e/f/15/0.5,
      d/t/10/0.5, e/t/10/0.5, f/t/10/0.5}
    \draw[fr] (\u) -- node[lb, pos=\p] {\c} (\v);
\end{tikzpicture}
\end{center}
La rete delle slide, con i nodi interni chiamati $a, \dots, f$. Le slide la riempiono un
cammino alla volta: il primo è $s \to a \to b \to e \to t$, che porta 4 unità, perché
l'arco $a \to b$ ha capacità 4 ed è il \kw{collo di bottiglia} del cammino.

\begin{teorema}[flusso massimo e taglio minimo]
Un \kw{taglio} $(S, T)$ divide i nodi in due parti, con $s \in S$ e $t \in T$; la sua
capacità è la somma delle capacità degli archi che vanno \emph{da} $S$ \emph{a} $T$. Il
valore del flusso massimo è uguale alla capacità minima di un taglio.
\end{teorema}

L'intuizione delle slide: aumentare la capacità di un arco collo di bottiglia aumenta il
flusso massimo, togliere un arco su un cammino critico lo diminuisce. Il taglio minimo è
proprio l'insieme dei colli di bottiglia: ogni unità di flusso deve attraversarlo.

\begin{esempio}[flusso massimo della rete]
Il flusso massimo vale \textbf{25}. Un taglio minimo è $S = \{s, c, f\}$: gli archi da $S$
a $T$ sono $s \to a$ (10), $s \to b$ (5) e $f \to t$ (10), in tutto 25. Gli archi
$b \to c$, $b \to f$ ed $e \to f$ vanno da $T$ a $S$ e non contano. Siccome flusso e taglio coincidono,
nessun flusso può superare 25 e nessun taglio può valere meno.
\end{esempio}

\begin{attenzione}[il taglio disegnato nelle slide]
La slide del taglio minimo evidenzia $S = \{s, b, c, f\}$ e scrive 5 sull'arco $b \to e$,
che nelle slide precedenti aveva capacità 8. Con la capacità originale quel taglio vale
$10 + 8 + 10 = 28$, non è minimo. Il taglio minimo è quello dell'esempio qui sopra; il
valore del flusso massimo, 25, non cambia (verificato con l'algoritmo di Edmonds--Karp e
con \texttt{networkx}, con entrambe le capacità).
\end{attenzione}

\begin{extra}[come si calcola: cammini aumentanti]
L'algoritmo di \kw{Ford--Fulkerson} parte dal flusso nullo e ripete: cerca un cammino da
$s$ a $t$ nella \kw{rete residua} (dove un arco con flusso $f$ e capacità $c$ ha ancora
$c - f$ in avanti e $f$ all'indietro, per poter ``disfare'' scelte precedenti) e ci spinge
il collo di bottiglia. Quando non ci sono più cammini, i nodi raggiungibili da $s$ formano
il taglio minimo. Scegliendo ogni volta il cammino più corto (\kw{Edmonds--Karp}) il costo
è $O(|V|\,|E|^2)$. Con capacità intere anche il flusso massimo trovato è intero
(\kw{teorema di integralità}): è ciò che rende corretta la riduzione del matching.
\end{extra}

\subsection{Matching massimo come flusso massimo}

Dati due insiemi $U$ e $V$ e le coppie ammesse fra loro, il \kw{matching di cardinalità
massima} usa più coppie possibili senza usare due volte lo stesso nodo. Si costruisce una
rete:
\begin{enumerate}
  \item una sorgente $s$ con un arco di capacità 1 verso ogni $u \in U$;
  \item un arco di capacità 1 da $u \in U$ a $v \in V$ per ogni coppia ammessa;
  \item un arco di capacità 1 da ogni $v \in V$ al pozzo $t$.
\end{enumerate}
Il flusso massimo è uguale al matching massimo. La capacità 1 sugli archi di $s$ e di $t$
fa sì che ogni nodo abbia al più un partner; la conservazione fa sì che il flusso che entra
in $u$ esca verso un solo $v$. Per l'integralità esiste un flusso massimo con valori 0 e 1,
e le coppie con flusso 1 sono il matching.

\begin{esempio}[le slide]
$U = \{u_1, u_2, u_3\}$, $V = \{v_1, v_2, v_3\}$, coppie ammesse $u_1v_1$, $u_1v_2$, $u_2v_2$,
$u_2v_3$, $u_3v_3$. Il flusso massimo vale 3: $\{u_1v_1, u_2v_2, u_3v_3\}$, un matching
perfetto. Le slide mostrano poi la versione con capacità: archi $s \to u_i$ di capacità 2,
e $v_2 \to t$ di capacità 2. Ora ogni $u$ può avere due partner e $v_2$ ne può avere due;
il flusso massimo vale 4, per esempio $\{u_1v_1, u_1v_2, u_2v_2, u_3v_3\}$. Non può valere
di più, perché gli archi verso $t$ portano in tutto $1 + 2 + 1 = 4$.
\end{esempio}

\subsection{Flusso a costo minimo}

Il flusso massimo guarda solo al volume. Il \kw{flusso a costo minimo} aggiunge l'economia:
ogni arco ha una capacità e un \kw{costo} per unità di flusso, e si vuole mandare da $s$ a
$t$ una quantità richiesta di flusso spendendo il meno possibile:
\[
  \text{costo} = \sum_{e \in E} c(e)\, f(e),
\]
dove qui $c(e)$ è il costo unitario dell'arco. Con i costi sugli archi $u \to v$ si
ottengono i problemi di \kw{assegnamento}: persone a compiti, revisori ad articoli, con il
costo che misura quanto è cattivo l'abbinamento.

\begin{esempio}[le slide]
Stessa rete del matching, con costi $u_1v_1 = 1$, $u_1v_2 = 5$, $u_2v_2 = 1$, $u_2v_3 = 1$,
$u_3v_3 = 5$, tutte le capacità 1. Mandare 3 unità significa trovare un matching perfetto, e
ce n'è uno solo: $u_3$ può andare solo con $v_3$, quindi $u_2$ va con $v_2$ e $u_1$ con
$v_1$. Costo $1 + 1 + 5 = 7$. Se l'arco $s \to u_2$ ha capacità 2, $u_2$ può prendere sia
$v_2$ sia $v_3$: il flusso di 3 unità diventa $u_1v_1$, $u_2v_2$, $u_2v_3$, con costo
$1 + 1 + 1 = 3$, e $u_3$ resta senza partner.
\end{esempio}

\begin{osservazione}[ottimizzare o stabilizzare]
Gli ultimi due esempi mostrano la differenza di prospettiva rispetto a Gale--Shapley.
Il flusso trova l'abbinamento migliore \emph{per un obiettivo globale} (più coppie, meno
costo), ma non si chiede se qualcuno vorrebbe deviare. Il matching stabile protegge dalle
deviazioni ma può essere lontano dall'ottimo globale. Quale serve dipende dal problema: in
un mercato dove i partecipanti possono accordarsi fuori dal sistema conta la stabilità; in
un'assegnazione decisa centralmente (compiti in un'azienda) conta l'obiettivo.
\end{osservazione}

% =====================================================================
\partelab{implementare e verificare gli algoritmi}
% =====================================================================

\textbf{Nota}: per questa lezione non c'è ancora un notebook del laboratorio. Questa parte
implementa gli algoritmi delle slide in Python e li verifica sugli esempi e su istanze
casuali, contro una ricerca esaustiva. Tutti i numeri riportati sono misurati. Quando
arriverà il notebook, la sezione si aggiornerà.

\subsection{Dittatura seriale e Top Trading Cycle}

\begin{lstlisting}
def serial_dictatorship(prefs, order):
    taken, out = set(), {}
    for a in order:
        item = next(x for x in prefs[a] if x not in taken)
        out[a] = item
        taken.add(item)
    return out

def top_trading_cycles(prefs, owner):
    """prefs: agente -> lista di oggetti; owner: oggetto -> agente."""
    agents, items, out = set(prefs), set(owner), {}
    while agents:
        top = {a: next(x for x in prefs[a] if x in items) for a in agents}
        a = min(agents)                      # da un agente qualunque...
        seen = []
        while a not in seen:                 # ...si seguono le frecce
            seen.append(a)
            a = owner[top[a]]
        cycle = seen[seen.index(a):]         # il primo nodo ripetuto chiude il ciclo
        for b in cycle:
            out[b] = top[b]
        agents -= set(cycle)
        items -= {top[b] for b in cycle}
    return out
\end{lstlisting}
Gli agenti e gli oggetti si riducono di almeno uno a ogni giro, e ogni giro costa
$O(n^2)$ nel caso peggiore: $O(n^3)$ in tutto, più che sufficiente per esempi didattici.
Sull'esempio delle slide \texttt{top\_trading\_cycles} restituisce
\texttt{\{A: 3, C: 1, B: 4, D: 2\}}; \texttt{serial\_dictatorship} con l'ordine
\texttt{"ABC"} restituisce \texttt{\{A: X, B: Z, C: Y\}}.

\paragraph{Verifiche.} Su 300 istanze casuali con 5 agenti, il risultato del TTC è sempre
individualmente razionale e Pareto efficiente (confronto con tutte le $5! = 120$
allocazioni). Su 200 istanze con 4 agenti, per ogni agente si sono provate tutte le
$4! = 24$ liste false: \textbf{nessuna} manipolazione migliora l'oggetto ricevuto, come
prevede la strategyproofness.

\subsection{Gale--Shapley e controllo della stabilità}

\begin{lstlisting}
from collections import deque

def gale_shapley(men, women):
    rank = {w: {m: i for i, m in enumerate(p)} for w, p in women.items()}
    nxt = {m: 0 for m in men}            # prossima donna da provare
    fiance = {}                          # donna -> uomo
    free = deque(men)
    while free:
        m = free.popleft()
        w = men[m][nxt[m]]
        nxt[m] += 1
        cur = fiance.get(w)
        if cur is None:
            fiance[w] = m
        elif rank[w][m] < rank[w][cur]:  # w preferisce m
            fiance[w] = m
            free.append(cur)
        else:
            free.append(m)
    return {m: w for w, m in fiance.items()}

def blocking_pairs(men, women, match):
    husband = {w: m for m, w in match.items()}
    out = []
    for m in men:
        for w in men[m]:                 # donne che m preferisce alla sua
            if w == match[m]:
                break
            if women[w].index(m) < women[w].index(husband[w]):
                out.append((m, w))
    return out
\end{lstlisting}
La tabella \texttt{rank} permette a una donna di confrontare due uomini in $O(1)$: senza,
ogni confronto costerebbe $O(n)$ e l'algoritmo $O(n^3)$. Per trovare \emph{tutti} i
matching stabili di un'istanza piccola basta provare le $n!$ permutazioni e tenere quelle
senza coppie bloccanti.

\paragraph{Verifiche contro la ricerca esaustiva.} Su 1000 istanze casuali con $n$ da 2 a
6 il risultato è sempre stabile, ogni uomo ha la migliore delle sue partner stabili e ogni
donna la peggiore dei suoi: M-ottimale e W-pessimale, come dice il teorema. Con $n = 6$ le
istanze hanno in media 2,08 matching stabili (da 1 a 8); in 69 casi su 200 ce n'è uno
solo. Sugli esempi delle slide si ritrovano il matching $\{(A,X),(B,Y),(C,Z)\}$ con 5
proposte e i due soli matching stabili della dimostrazione di Roth.

\paragraph{Chi può manipolare.} Su 300 istanze casuali con $n = 4$, provando per ogni
agente tutte le 24 liste false: con gli uomini che propongono, \textbf{nessun uomo}
migliora mai mentendo; in 12 istanze su 300 (4\%) almeno una donna sì. È il teorema di
Roth visto dal vivo: la manipolazione esiste, ma solo dal lato che riceve le proposte, e
su istanze casuali è rara.

\paragraph{Quante proposte servono davvero.} Il limite è $n^2$, ma con preferenze casuali
le proposte sono molte meno, circa $n \ln n$ (media su 5 istanze):
\begin{center}
\renewcommand{\arraystretch}{1.2}
\begin{tabular}{rrrrr}
\toprule
$n$ & Proposte & $n \ln n$ & $n^2$ & Tempo \\
\midrule
100  & 459     & 461     & 10\,000     & 0,5 ms \\
500  & 3\,208  & 3\,107  & 250\,000    & 15 ms \\
1000 & 6\,669  & 6\,908  & 1\,000\,000 & 67 ms \\
2000 & 14\,975 & 15\,202 & 4\,000\,000 & 275 ms \\
\bottomrule
\end{tabular}
\end{center}
I tempi crescono più delle proposte perché costruire \texttt{rank} legge tutte le $n^2$
preferenze: è il costo di leggere l'input, inevitabile.

\subsection{Capacità}

Basta che ogni scuola tenga i suoi migliori $c_i$ e rimandi indietro il peggiore quando
supera la capacità:
\begin{lstlisting}
def deferred_acceptance_cap(students, schools, cap):
    rank = {h: {s: i for i, s in enumerate(p)} for h, p in schools.items()}
    nxt = {s: 0 for s in students}
    held = {h: [] for h in schools}
    free = deque(students)
    while free:
        s = free.popleft()
        if nxt[s] >= len(students[s]):   # lista finita: resta senza scuola
            continue
        h = students[s][nxt[s]]
        nxt[s] += 1
        held[h].append(s)
        held[h].sort(key=lambda x: rank[h][x])
        if len(held[h]) > cap[h]:
            free.append(held[h].pop())   # rimanda indietro il peggiore
    return held
\end{lstlisting}
Sull'esempio delle slide restituisce $X = \{A, E\}$, $Y = \{B, C, D\}$. Il controllo della
lista finita gestisce anche le liste incomplete.

\subsection{Flussi con NetworkX}

Il flusso massimo e il taglio minimo della rete delle slide:
\begin{lstlisting}
import networkx as nx
G = nx.DiGraph()
for u, v, c in [("s","a",10), ("s","b",5), ("s","c",15), ("a","b",4),
                ("b","c",4), ("a","d",9), ("a","e",15), ("b","e",8),
                ("b","f",6), ("c","f",13), ("d","e",15), ("e","f",15),
                ("d","t",10), ("e","t",10), ("f","t",10)]:
    G.add_edge(u, v, capacity=c)
valore, (S, T) = nx.minimum_cut(G, "s", "t")   # 25, S = {s, c, f}
\end{lstlisting}
Un'implementazione diretta di Edmonds--Karp (ricerca in ampiezza nella rete residua) dà lo
stesso valore 25 e lo stesso taglio. Per il flusso a costo minimo l'attributo
\texttt{weight} è il costo unitario:
\begin{lstlisting}
H = nx.DiGraph()
for u in ["u1", "u2", "u3"]:
    H.add_edge("s", u, capacity=1, weight=0)
for (u, v), costo in {("u1","v1"): 1, ("u1","v2"): 5, ("u2","v2"): 1,
                      ("u2","v3"): 1, ("u3","v3"): 5}.items():
    H.add_edge(u, v, capacity=1, weight=costo)
for v in ["v1", "v2", "v3"]:
    H.add_edge(v, "t", capacity=1, weight=0)
F = nx.max_flow_min_cost(H, "s", "t")
nx.cost_of_flow(H, F)                          # 7
\end{lstlisting}
Con \texttt{capacity=2} sull'arco $s \to u_2$ il costo scende a 3, con il flusso su
$u_1v_1$, $u_2v_2$, $u_2v_3$: i due casi delle slide. Per il solo matching massimo in un
grafo bipartito c'è anche \texttt{nx.bipartite.maximum\_matching} (Hopcroft--Karp), che
non richiede di costruire la rete.

% ---------------------------------------------------------------------
\begin{riepilogo}
\begin{itemize}
  \item \textbf{Tre famiglie}: preferenze da un lato (efficienza, incentivi), da due lati
        (stabilità), ottimizzazione (flussi).
  \item \textbf{Dittatura seriale}: in ordine, ognuno prende il migliore rimasto. Pareto
        efficiente e strategyproof, non massimizza il benessere, equa solo in valore atteso
        con ordine casuale. Round robin: più equo, manipolabile.
  \item \textbf{TTC}: ognuno punta al preferito, l'oggetto al proprietario; si eseguono i
        cicli. Individualmente razionale, unica allocazione nel core, strategyproof.
  \item \textbf{Matrimonio stabile}: nessuna coppia bloccante. Gale--Shapley trova sempre
        un matching stabile in al più $n^2$ proposte; chi propone ottiene il suo ottimo
        stabile, l'altro lato il suo peggiore.
  \item \textbf{Roth}: nessun meccanismo è stabile e strategyproof insieme; l'accettazione
        differita è strategyproof solo per chi propone.
  \item \textbf{Estensioni}: $n \ne m$, liste incomplete, pareggi; capacità (studenti e
        scuole) con la stessa accettazione differita.
  \item \textbf{Flussi}: flusso massimo = taglio minimo; il matching massimo è un flusso
        con capacità 1; il flusso a costo minimo risolve l'assegnamento. Nell'esempio delle
        slide il flusso massimo è 25, e il taglio disegnato (28 con la capacità originale)
        non è quello minimo.
  \item \textbf{Lab}: tutti gli esempi verificati; su istanze casuali Gale--Shapley è
        M-ottimale, nessun proponente manipola, $\approx n\ln n$ proposte invece di $n^2$.
\end{itemize}
\end{riepilogo}
